\section{Introduction}
\label{sec:introduction}

Long-running language-model agents may rely on memories formed under conditions
that no longer hold. Stored observations and reflections support persistent
simulated agents \citep{park2023generative}, verbal feedback can guide later
trials \citep{shinn2023reflexion}, and trajectory-level lessons can transfer
across tasks \citep{zhao2023expel}. Yet an experience that was useful when
recorded can become stale as operating conditions change. This risk is
especially consequential when the memory recommends an operational strategy and
its outcome is observable only after delayed execution. The resulting question
is not simply what an agent should remember, but when it should continue to
trust its past.

Existing research addresses parts of this problem without resolving the
operational decision studied here. Long-term memory benchmarks test temporal
reasoning and knowledge updates
\citep{maharana2024locomo,wu2025longmemeval}. Concept-drift research distinguishes
change detection from the adaptation that follows \citep{gama2014survey}, while
nonstationary inventory research develops control under changing or partially
observed demand \citep{treharne2002adaptive,bayraktar2010inventory}. Together,
these literatures address persistence, change detection, and adaptive control,
but do not establish whether an explicit memory lifecycle reduces downstream
cost under delayed feedback. A detector can indicate that public observations
changed. It cannot by itself establish that a
retrieved experience is invalid or that suspending it will reduce downstream
cost.

We study this question in a synthetic single-item lost-sales inventory system,
chosen as a controlled testbed rather than a model of enterprise deployment.
ShiftMem treats a change in public context as a prompt for review rather than
proof of invalidity, and uses delayed outcomes to update memory confidence and
state. Page--Hinkley signals provide the change trigger
\citep{page1954continuous}. The LLM is confined to bounded strategy revisions,
while a deterministic controller issues daily orders. We compare ShiftMem with
a lexical retrieval baseline under paired demand and supply streams. This design
compares the end-to-end systems. It does not isolate the effect of any lifecycle
component.

The prespecified evaluation covered two provider-hosted model deployments,
eight held-out scenarios, and 160 completed
deployment--method--scenario--seed runs. Its dependence-limited primary
decision rule did not support a ShiftMem advantage, and the mean point estimate
was unfavorable. Because the 70 model-specific pairs shared 35 environment
clusters, we also report a post-hoc clustered sensitivity, which retained the
same unfavorable direction. Point estimates varied descriptively across test
regimes and deployments without establishing generalizable effect modification.
The lexical retrieval baseline always ran first, leaving method order and provider
drift as potential confounders.

\paragraph{Contributions.}
This study makes three bounded contributions:
\begin{itemize}
  \item It develops a controlled architecture that separates bounded LLM
  strategy review from deterministic daily execution, allowing memory-conditioned
  revision to be evaluated without changing the inventory controller.
  \item It provides a content-addressed evidence package with network-free
  reproducible integrity checks and deterministic post-outcome analyses. The package records
  provider and structured-output failures, fallback behavior, reservation
  journals, and source hashes.
  \item It reports the negative primary result with its uncertainty and
  dependence-aware sensitivity analyses, together with descriptive deployment
  heterogeneity and post-hoc reliability analyses.
\end{itemize}
